\section{The actual original Cartier ratio section and intrinsic ideal potentials}

\begin{definition}[Canonical original generic Cartier ratio section]
\label{def:actual-cartier-ratio-generic-section}
\lean{EnvelopeProperty.Berkovich.RelativeSchemeValuedPoint.genericCartierRatioSection}
\leanok
\uses{def:actual-cartier-generic-inclusion-frame,cor:actual-dual-cartier-fiber-section-model-potential,thm:actual-model-line-tensor-norm}
For two actual Cartier ideals $I,J$ on an original proper model,
and nonzero intrinsic original ideal norm of $I$, take the genuine
canonical generic unit section of the original ideal line $I$
and the actual original fiber section of the canonical dual
transpose of the original inclusion of $J$. Their true explicit
pure tensor is an actual section on the same original generic
fiber of the original line $I\otimes J^\vee$.
\end{definition}

\begin{theorem}[Actual Cartier ratio model norm and logarithm]
\label{thm:actual-cartier-ratio-generic-model-norm}
\lean{EnvelopeProperty.Berkovich.RelativeSchemeValuedPoint.modelLineNorm_genericCartierRatioSection}
\lean{EnvelopeProperty.Berkovich.RelativeSchemeValuedPoint.modelLineNorm_genericCartierRatioSection_log}
\leanok
\uses{def:actual-cartier-ratio-generic-section,lem:actual-cartier-generic-unit-section-model-norm,cor:actual-dual-cartier-fiber-section-model-potential}
The genuine model norm of this actual original ratio section is
$\|J\|_y/\|I\|_y$, using the intrinsic original ideal norms. Its
original logarithmic norm is
$\log\|J\|_y-\log\|I\|_y$.
\end{theorem}
\begin{proof}
The actual explicit tensor norm theorem gives the product of the
two genuine original factor norms. The canonical original generic
unit section of $I$ has model norm reciprocal to the actual
Cartier inclusion norm. The actual original fiber section of the
canonical dual inclusion transpose for $J$ has exactly the actual
original Cartier inclusion norm of $J$. Each inclusion norm is
its corresponding intrinsic original model ideal norm. Taking
the actual extended-real logarithm of this exact norm identity,
using its original multiplication and inverse laws, gives the
asserted difference with the indicated order. This order records
the distinction between the logarithmic norm and the usual
negative-logarithmic model weight. No tensor, dual metric, or
rational trivialization certificate has been assumed.
\end{proof}

\begin{corollary}[Exact original ratio model weight]
\label{cor:actual-cartier-ratio-model-weight}
\lean{EnvelopeProperty.Berkovich.RelativeSchemeValuedPoint.modelLineNorm_genericCartierRatioSection_weight}
\leanok
\uses{thm:actual-cartier-ratio-generic-model-norm,scheme:native-model-ideal-function}
The negative original logarithmic norm of this actual ratio
section is exactly
$\log\|I\|_y-\log\|J\|_y$, the original ideal potential difference.
\end{corollary}
\begin{proof}
The intrinsic ideal norm of $I$ is nonzero by hypothesis and at
most one by the actual valuation-ring bound. Its logarithm is
therefore a finite original extended real. The valid extended-real
negation-of-difference identity consequently applies even if the
original $J$ norm vanishes. Negating the preceding exact logarithmic
norm formula yields the asserted potential difference with the
proper order. No extra finiteness certificate is needed.
\end{proof}
